%%%%%%%%%%%%%%%%%%%%%%%%%%%%%%%%%%%%%%%%%%%%%%%%%%%%%%%
% % Lines starting with % are comments, which are ignored.
% % This is a handy way of indicating the date and version of
% % your document, to wit:
% %
% % LaTeX sample file
% % Modified March, 2002
% %
%%%%%%%%%%%%%%%%%%%%%%%%%%%%%%%%%%%%%%%%%%%%%%%%%%%%%%%
\DocumentMetadata{
  % PDF standards: ONE key whose value is a comma-separated list --
  % UA-2 (accessibility) plus A-4f (archival; the archival companion of
  % UA-2). Do NOT use separate pdfstandard=... keys: with ordinary
  % key-value parsing only the last one takes effect.
  pdfstandard={UA-2,A-4f},
  lang=en-US,
  tagging=on,
  tagging-setup={
    math/setup=mathml-SE
  }
}

\documentclass{ltx-talk}
% lecture-macros.tex -- shared package setup for lecture-template.tex
% (ltx-talk slides / handouts).
%
% NOTE ON PATHS: every template in this repository is stored one level up
% from where it is compiled (see README.md). All relative \input paths
% below are resolved against the folder the document is compiled in,
% which sits one directory above this file.

% --- Tables, bookmarks, strikethrough ------------------------------------
% booktabs: professional table rules (\toprule, \midrule, \bottomrule)
% bookmark: PDF outline entries (PDF bookmarks)
% ulem[normalem]: \sout strikethrough that survives math mode
\usepackage{booktabs}
\usepackage{bookmark}
\usepackage[normalem]{ulem}
\usepackage{multicol}

% --- Math ------------------------------------------------------------
% amsmath: equation environments and numbering
% mathrsfs: \mathfrak  |  stmaryrd: extra symbols
\usepackage{amsmath}
\usepackage{mathrsfs}
\usepackage{stmaryrd}

% --- Layout -------------------------------------------------------------
% float: [H] float placement
% fontspec: Unicode/OpenType fonts (compile with XeLaTeX or LuaLaTeX)
\usepackage{float}
\usepackage{fontspec}

% --- Colors ---------------------------------------------------------------
% dvipsnames: extended named color palette (the Accessible* colors below
% are built on it);
% table: \rowcolors / \cellcolor for shading tables.
% NOTE: xcolor may only be loaded ONCE per document -- a second
% \usepackage{xcolor} with different options is a fatal "Option clash"
% error, so ALL of its options are listed here.
\usepackage[dvipsnames,table]{xcolor}

% --- Multimedia -----------------------------------------------------------
% animate: embeds audio/video (and legacy Flash) media in a frame.
\usepackage{animate}

% --- Bibliography ----------------------------------------------------------
% biblatex (biber backend): citations for lecture notes.
\usepackage[
  backend=biber,
  style=alphabetic,
  url=true,
  doi=true,
  isbn=true,
  backref=true
]{biblatex}

% --- Colored boxes ----------------------------------------------------------
% talkthemetcolorbox: the colored-box theme for ltx-talk slides
% (shadow/rounded/shaded blocks, showtitle, titlepage styles).
\usepackage[
  shadow,
  rounded,
  shaded,
  blocks,
  showtitle,
  titlepage
]{talkthemetcolorbox}

% --- "theindex" helper environment ------------------------------------------
% Provides a lightweight theindex environment (an itemize) and the
% \@idxitem/\subitem/\subsubitem/\indexspace macros that LaTeX's standard
% theindex machinery expects, so hand-written index pages indent correctly.
\makeatletter
\newenvironment{theindex}
  {\begin{itemize}}
  {\end{itemize}}
\providecommand\@idxitem{\par\hangindent 40\p@}
\providecommand\subitem{\@idxitem \hspace*{20\p@}}
\providecommand\subsubitem{\@idxitem \hspace*{30\p@}}
\providecommand\indexspace{\par \vskip 10\p@ \@plus5\p@ \@minus3\p@\relax}
\makeatother

% --- Index -----------------------------------------------------------------
% imakeidx: builds the printed index.
% The redefinition below appends "|hyperpage" to every index entry so that,
% in the PDF, each index entry is a link to its page.
\usepackage{imakeidx}
\makeatletter
\let\imki@orig@wrindexentrysplit\imki@wrindexentrysplit
\renewcommand\imki@wrindexentrysplit[3]{%
  \imki@orig@wrindexentrysplit{#1}{#2|hyperpage}{#3}%
}

\makeatletter
% --- Section titles: capture + one bookmark per logical section -------------
% \section is redefined so that:
%   1. the section title is captured into the plain-TeX register
%      \mytalksectitletoks (reprinted verbatim on the auto divider frame,
%      see the cmd/section/after hook below -- plain-TeX storage is
%      catcode-safe, unlike expanding the title in a hook), and
%   2. exactly ONE depth-1 PDF bookmark is created per \section call
%      (hyperref's own bookmark is suppressed via \Hy@writebookmark).
\DeclareCommandCopy{\OriginalSectionForBookmarks}{\section}
\newtoks\mytalksectitletoks

\RenewDocumentCommand{\section}{s o m}{%
  \begingroup
    \def\Hy@writebookmark##1##2##3##4##5{}%
    \global\mytalksectitletoks={#3}%
    \IfBooleanTF{#1}
      {\phantomsection\IfNoValueTF{#2}{\OriginalSectionForBookmarks*{#3}}{\OriginalSectionForBookmarks*[#2]{#3}}}%
      {\IfNoValueTF{#2}{\OriginalSectionForBookmarks{#3}}{\OriginalSectionForBookmarks[#2]{#3}}}%
  \endgroup
  \IfBooleanTF{#1}
    {\IfNoValueTF{#2}{\pdfbookmark[1]{#3}{section-star.\thepage.\number\inputlineno}}{\pdfbookmark[1]{#2}{section-star.\thepage.\number\inputlineno}}}%
    {\IfNoValueTF{#2}{\pdfbookmark[1]{Section \thesection: #3}{section.\thesection}}{\pdfbookmark[1]{Section \thesection: #2}{section.\thesection}}}%
}
\makeatother

% --- One PDF bookmark per logical frame (NOT per overlay) ---------------------
% ltx-talk writes one PDF page per overlay, so a naive frametitle hook would
% register a new bookmark for every overlay of the same frame. The
% \g_slideutils_frame_bookmark_int register remembers the last frame that
% got a bookmark, so subsequent overlays of the same frame are skipped.
% This is also what keeps BeamerPresenter's cumulative annotations working:
% one logical frame must stay one logical unit in the PDF outline.
\ExplSyntaxOn
\int_new:N \g_slideutils_frame_bookmark_int

\AddToHook{cmd/frametitle/after}{%
  \tl_if_empty:NF \g__talk_frame_title_tl
    {
      \ifnum\value{frame}=\g_slideutils_frame_bookmark_int\relax
      \else
        \int_gset:Nn \g_slideutils_frame_bookmark_int { \value{frame} }%
        \pdfbookmark[2]{\tl_use:N \g__talk_frame_title_tl}
          {frame.\int_use:N \g_slideutils_frame_bookmark_int}
      \fi
    }
}
\ExplSyntaxOff
\makeatother

% --- Printed index + theme for the colored boxes -------------------------------
% \makeindex[...]: two-column index titled "Index", listed in the TOC.
% The two EditInstance calls style the standard header/footer instances
% of talkthemetcolorbox (cyan background; footer order title/author/
% date/framenumber).
\makeindex[columns=2, title=Index, intoc]\EditInstance{header}{std}{
  background-color = cyan,
  color            = black,
  height           = 0.75cm
}
\EditInstance{footer}{std}{
  background-color = cyan,
  element-order = {title, author, date, framenumber},
}

% --- Auto section divider frame ---------------------------------------------------
% After each \section, insert a centered, large, structure-colored frame
% that prints the section title (from \mytalksectitletoks).
\AddToHook{cmd/section/after}{%
  \begin{frame}
    \centering
    \color{structure}
    \LARGE
    \the\mytalksectitletoks
  \end{frame}%
}

% --- Tagging sockets (tagged PDF structure) ---------------------------------------
% ltx-talk exposes tagging "sockets" where the PDF structure can attach
% section titles; register one and leave it unassigned (none) so the
% class's own tagging keeps control.
\ExplSyntaxOn
\AtBeginDocument{
  \NewTaggingSocketPlug { talk / sec / title } { none } { }
  \AssignTaggingSocketPlug { talk / sec / title } { none }
}
\ExplSyntaxOff

% --- Accessibility roles ----------------------------------------------------------------
% Map \frametitle to the H1 heading role in the tagged PDF, so screen
% readers treat slide titles as headings.
\tagpdfsetup{role / new-tag = frametitle / H1}

\makeatletter
% --- Page label = frame number -------------------------------------------------------
% Redefine what feeds the PDF page label to track the frame counter,
% not the physical page counter. Replace \theframenumber below with
% whatever internal macro/property ltx-talk actually uses to produce
% the "framenumber" footer element — confirm via ltx-talk-frame.dtx.
\renewcommand\thepage{\arabic{frame}}
\makeatother

\addbibresource{refs.bib}

% --- "Problem Time!" slides ------------------------------------------------------------------
% \ProblemSlide{problem}{solution} emits two frames: the problem, then
% its solution after a \pause, each numbered by the problem counter.
\newcounter{problemcounter}
\setcounter{problemcounter}{0}
\newcommand{\ProblemSlide}[2]{
  \stepcounter{problemcounter}
  \begin{frame}
    \frametitle{Problem Time!}
    Problem \theproblemcounter: #1
  \end{frame}

  \begin{frame}
    \frametitle{Problem Time!}
    Solution to Problem \theproblemcounter: \pause

    #2
  \end{frame}
}
% \highlight[<colour>]{<stuff>}
% from https://tex.stackexchange.com/questions/33401/a-version-of-colorbox-that-works-inside-math-environments
\newcommand{\highlight}[2][yellow]{\mathchoice%
  {\colorbox{#1}{$\displaystyle#2$}}%
  {\colorbox{#1}{$\textstyle#2$}}%
  {\colorbox{#1}{$\scriptstyle#2$}}%
  {\colorbox{#1}{$\scriptscriptstyle#2$}}}%

% --- Fonts ------------------------------------------------------------------------------
% Lato (text) + LeteSansMath (math): open, accessible typefaces.
\setmainfont{Lato}[
  Ligatures=TeX,
  Numbers={Lining,Monospaced}
]
\setsansfont{Lato}[
  Ligatures=TeX,
  Numbers={Lining,Monospaced}
]
\setmathfont{LeteSansMath.otf}\usepackage{emoji}

% --- Accessible colors ---------------------------------------------------------------------
% WCAG-contrast-safe variants of red/green for text on white:
% red!80!black and green!40!black both clear the 4.5:1 requirement.
% AltYellow: soft background for highlight boxes.
\colorlet{AccessibleRed}{red!80!black}
\colorlet{AccessibleGreen}{green!40!black}
\definecolor{AltYellow}{RGB}{255,245,170}

% --- Links ---------------------------------------------------------------------------------
% hyperref + link colors: every link uses AccessibleGreen (7.2:1 on
% white), so internal, citation, and URL links all pass WCAG 1.4.3.
\usepackage{hyperref}
\hypersetup{
  colorlinks=true,
  linkcolor=AccessibleGreen,
  urlcolor=AccessibleGreen,
  citecolor=AccessibleGreen,
}

% --- Shared course macros ---------------------------------------------------------------------
% Root macros.tex: math shorthands used across all course documents
% (\RR, \NN, \ZZ, \QQ, \grad, \curl, \tr, \der, \pder, ...).
\newcommand{\RR}{\mathbb{R}}
\newcommand{\NN}{\mathbb{N}}
\newcommand{\ZZ}{\mathbb{Z}}
\newcommand{\QQ}{\mathbb{Q}}
\newcommand{\vv}{\vec{v}}
\newcommand{\ww}{\vec{w}}
\newcommand{\xx}{\vec{x}}
\newcommand{\yy}{\vec{y}}
\newcommand{\zz}{\vec{z}}
\newcommand{\uu}{\vec{u}}
\newcommand{\basis}[1]{\vec{e_{#1}}}
\newcommand{\id}{\mathrm{id}}
\newcommand{\CC}{\mathbb{C}}
\renewcommand{\div}{\operatorname{div}}
\newcommand{\grad}{\operatorname{grad}}
\newcommand{\curl}{\operatorname{curl}}
\newcommand{\tr}{\operatorname{tr}}
\newcommand{\der}[2]{\dfrac{d #1}{d #2}}
\newcommand{\pder}[3][]{\dfrac{\partial^{#1} #2}{\partial #3}}


% % Title and author(s)
%%%%%%%%%%%%%%%%%%%%%%%%%%%%%%%%%%%%%%%%%%%%%%%%%%%%%%%
\title{Lecture 11 - The Chain Rule}
\author{Garrett Credi}
\date{10/15/2026}
%%%%%%%%%%%%%%%%%%%%%%%%%%%%%%%%%%%%%%%%%%%%%%%%%%%%%%%
% %
% % The next command allows your in import encapsulated
% % postscript files, .epsf or .eps files, which
% % contain vector graphic image data.
% %
%%%%%%%%%%%%%%%%%%%%%%%%%%%%%%%%%%%%%%%%%%%%%%%%%%%%%%%
% %   The first thanks indicates your affiliation
% %
% %  Just the name here.
% %
% % Your mailing address goes at the end.
% %
% % \thanks is also how you indicate grant support
% %
%%%%%%%%%%%%%%%%%%%%%%%%%%%%%%%%%%%%%%%%%%%%%%%%%%%%%%%

%%%%%%%%%%%%%%%%%%%%%%%%%%%%%%%%%%%%%%%%%%%%%%%%%%%%%%%
% $ The next command ensures the margins are 1 inch on
% % all sides.
% %
%%%%%%%%%%%%%%%%%%%%%%%%%%%%%%%%%%%%%%%%%%%%%%%%%%%%%%%
\begin{document}
\begin{frame}
  \maketitle 
\end{frame}

\begin{frame}
  \frametitle{Quiz Time}

  Let's get into neat rows, spread out, and slay this quiz! :)
\end{frame}

\section{The Chain Rule}

\begin{frame}
  \frametitle{The Road So Far}
    Recall we're trying to build a \textit{synthetic} approach to studying derivatives: \pause

    Atomic Rules: \pause
  \begin{multicols}{2}
    \begin{itemize}
      \item Power Rule: $\dfrac{d}{dx} x^n$ 
      \item Trigonometric Rules: $\dfrac{d}{dx} \sin(x)$ etc.
      \item Exponential/Logarithm Rules: $\dfrac{d}{dx} e^x$, $\dfrac{d}{dx} \ln(x)$
    \end{itemize}
  \end{multicols}

  Combination Rules: \pause
  \begin{multicols}{2}
    \begin{itemize}
      \item Sum/Difference Rule: $\dfrac{d}{dx} (f(x) \pm g(x))$ 
      \item Product Rule: $\dfrac{d}{dx} f(x)g(x)$ 
      \item Chain Rule: $\dfrac{d}{dx} g(f(x))$ 
      \item Constant Multiple Rule: $\dfrac{d}{dx} c \cdot f(x)$
      \item Quotient Rule: $\dfrac{d}{dx} \dfrac{f(x)}{g(x)}$ 
      \item Inverse Function Rule: $\dfrac{d}{dx} f^{-1}(x)$
    \end{itemize}
  \end{multicols}
\end{frame}

\begin{frame}
  \frametitle{What is the Chain Rule? I}

  We're spending today building up the chain rule, the tool that will help us compute derivatives of the form: \pause

  \begin{equation}
    \dfrac{d}{dx} ( g (f(x)))
  \end{equation}

  But what kinds of functions will this actually help us compute derivative of? \pause

  Sum/Difference Rule and Product/Quotient rule have visible signifiers that help us determine when to use them: \pause
  \vspace{1in}

  But what do functions of the form $g(f(x))$ even look like?
\end{frame}

\begin{frame}
  \frametitle{What is the Chain Rule? II}

  Recall our ``basic" functions: \pause

  \begin{multicols}{3}
    \begin{itemize}
      \item $x^n$
      \item $c^x$
      \item $\sin(x), \cos(x)$, etc
      \item $|x|$
      \item $\log_c(x)$
    \end{itemize} \pause

  \end{multicols}

  When we \textit{compose} functions (do the $g(f(x))$ thing), what does this actually do? \pause

  We make one function's output another function's input instead of just inputting $x$! 
\end{frame}

\begin{frame}
  \frametitle{What is the Chain Rule? III}

  So, before we even \textit{use} chain rule, we should first come up with a strategy for: \pause
  \begin{enumerate}
    \item Recognizing when functions are composite functions \pause
    \item Recognizing which function is the inner function (the $f(x)$ in the $g(f(x))$) and which function is the outer function (the $g(x)$ in the $g(f(x))$) \pause
  \end{enumerate}

  Strategy: \pause
  \begin{enumerate}
    \item Break up the complicated function into the addition, subtraction, product, and quotient of simpler functions (i.e. do sum/difference and product/quotient rules first) \pause
    \item In each of the smaller pieces, identify where the variables actually are. \pause
    \item At each occurrence of a variable, find the biggest/longest expression that it lives in - this is the innermost function. \pause
    \item Treat that whole expression as a new variable, and repeat from step 1.
  \end{enumerate}
\end{frame}

\begin{frame}
  \frametitle{Identifying Composites Example}

  \vspace{\stretch{0.2}}
  For instance, let's take $h(t) = 3 \sin^2(4t^2 + 1) \cdot \cos (4t^2 + 1) \cdot 8t$
  \vspace{\stretch{1}}
\end{frame}

\begin{frame}
  \frametitle{Deriving Chain Rule I}

  Now let's assume we've already identified some function of interest as a composite $g(f(x))$, \pause
  we'd like to find an expression of the derivative $\dfrac{d}{dx} g(f(x))$ hopefully in terms of derivatives involving $f(x)$ and $g(x)$. \pause

  Strategy: Linearize once again! \pause Let's linearize around some $x = x_0$. \pause

  We start by linearizing \textit{only} $f(x)$, \pause since our input to $g(x)$ is changed by $f$ and thus might end up close to a point different than $x = x_0$. \pause

  We know that $f(x) \approx f(x_0) + f'(x_0)(x - x_0)$, \pause so we end up wanting to further linearize:
  \begin{equation*}
    g(f(x_0) + f'(x_0)(x - x_0))
  \end{equation*}
\end{frame}

\begin{frame}
  \frametitle{Deriving Chain Rule II}

  Note: when we linearize $f(x)$ near $x = x_0$, we mean that for values $x \approx x_0$ we'll know that $f(x) \approx f(x_0) + f'(x_0)(x - x_0)$ \pause
  i.e. our input value being close to something implies our output value is close to the related thing. \pause

  For
  \begin{equation*}
    g(f(x_0) + f'(x_0)(x - x_0))
  \end{equation*}
  what value are we inputting into the function $g$? \pause

  When $x \approx x_0$, what is $f(x_0) + f'(x_0)(x - x_0)$ close to? \pause

  Since $f(x_0) + f'(x_0) (x - x_0) \approx f(x_0)$, we will linearize $g(x)$ around $f(x_0)$ instead!
\end{frame}

\begin{frame}
  \frametitle{Deriving Chain Rule III}

  \begin{align}
    g(f(x)) &\approx g(f(x_0) + f'(x_0)( x-x_0)) \quad \text{linearize $g$ around $f(x_0)$} \\ \pause
            &\approx g(f(x_0)) + g'(f(x_0)) ( (f(x_0) + f'(x_0)(x -x_0)) - f(x_0)) \\ \pause
            &= g(f(x_0)) + g'(f(x_0))(f'(x_0)(x - x_0) + f(x_0) - f(x_0)) \\ \pause
            &= g(f(x_0)) + g'(f(x_0))f'(x_0)(x - x_0)
  \end{align} \pause
\end{frame}
\begin{frame}
  \frametitle{Deriving Chain Rule IV}
  The derivative is the slope of the best linear approximation $\implies$ \pause

  \begin{alertblock}{The Chain Rule}\index{derivative!chain rule}
    \begin{equation}
      \dfrac{d}{dx} g(f(x)) = g'(f(x)) f'(x)
    \end{equation}
  \end{alertblock} \pause

  Or, written differently

  \begin{alertblock}{The Chain Rule (Leibniz's Version)}
    \begin{equation}
      \dfrac{d (g \circ f)}{dx} \mid_{x = x_0} = \dfrac{dg}{dx} \mid_{x = f(x_0)} \cdot \dfrac{df}{dx} \mid_{x = x_0}
    \end{equation}
  \end{alertblock}
\end{frame}

\begin{frame}
  \frametitle{Chain Rule Example I}
  \vspace{\stretch{0.2}}
  Let's compute the derivative of that function from before:
  \begin{equation}
    h(t) = 3 \sin^2(4t^2 + 1) \cdot \cos(4t + 1) \cdot 8t
  \end{equation}

  \vspace{\stretch{1}}
\end{frame}

\begin{frame}
  \frametitle{Chain Rule Example II}
  We can also derive the quotient rule in terms of the chain rule since
  \begin{equation}
    \dfrac{f(x)}{g(x)} = f(x) \cdot \dfrac{1}{g(x)} = f(x) \cdot (\dfrac{1}{x} \circ g)(x)
  \end{equation}
  i.e. $\dfrac{1}{g(x)}$ is a composite function! You first do $g(x)$ and then do the function $\dfrac{1}{x}$. \pause

  Then: 

  \begin{align}
    \dfrac{d}{dx} \dfrac{f(x)}{g(x)} &= \left(\dfrac{d}{dx} f(x)\right) \cdot \dfrac{1}{g(x)} + f(x) \cdot \left(\dfrac{d}{dx} \dfrac{1}{g(x)}\right) \\ \pause
                                     &= \dfrac{f'(x)}{g(x)} + f(x) \cdot \left(\left(\dfrac{d}{du} \dfrac{1}{u}\right)\mid_{u=g(x)} \cdot \left(\dfrac{d}{dx} g(x)\right)\right) \\ \pause
                                     &= \dfrac{f'(x)}{g(x)} + f(x) \cdot \dfrac{-1}{g(x)^2} \cdot g'(x) \pause
                                     = \dfrac{f'(x) g(x)}{g(x)^2)} - \dfrac{f(x)g'(x)}{g(x)^2} \\ \pause
                                     &= \dfrac{f'(x) g(x) - f(x) g'(x)}{g(x)^2}
  \end{align}
\end{frame}

\begin{frame}
  \frametitle{Chain Rule Example III}
  We can also frame this more in terms of real world applications, too long for a slide but it's in the worksheet.
\end{frame}

\begin{frame}
  \frametitle{Worksheet Time}

  Let's get into groups and work on the worksheet!
\end{frame}

\begin{frame}
  \frametitle{End of Class}

  That's all folks!
\end{frame}
\end{document}
